\documentclass[12pt]{article}

\usepackage[T1]{fontenc}
\usepackage[utf8]{inputenc}
\usepackage{mathpazo}
\usepackage{amsmath,amssymb,amsthm}
\usepackage{geometry}
\usepackage{microtype}
\usepackage{setspace}
\usepackage{titlesec}
\usepackage{hyperref}
\usepackage{booktabs}
\usepackage{array}
\usepackage{tabularx}
\usepackage{enumitem}
\usepackage{xcolor}
\usepackage{multirow}

\geometry{
  paper=a4paper,
  top=1.25in,
  bottom=1.25in,
  left=1.4in,
  right=1.4in
}

\onehalfspacing

\hypersetup{
  colorlinks=true,
  linkcolor=black,
  citecolor=black,
  urlcolor=black
}

\titleformat{\section}{\large\bfseries}{\thesection.}{0.6em}{}
\titleformat{\subsection}{\normalsize\bfseries}{\thesubsection.}{0.5em}{}
\titleformat{\subsubsection}{\normalsize\itshape}{\thesubsubsection.}{0.5em}{}

% Theorem environments
\newtheoremstyle{rthm}
  {10pt}{10pt}{\itshape}{0pt}{\bfseries}{.}{0.5em}{}
\theoremstyle{rthm}
\newtheorem{theorem}{Theorem}[section]
\newtheorem{corollary}[theorem]{Corollary}
\newtheorem{proposition}[theorem]{Proposition}
\newtheorem{lemma}[theorem]{Lemma}

\newtheoremstyle{rdef}
  {8pt}{8pt}{\normalfont}{0pt}{\bfseries}{.}{0.5em}{}
\theoremstyle{rdef}
\newtheorem{definition}[theorem]{Definition}
\newtheorem{assumption}[theorem]{Assumption}
\newtheorem{example}[theorem]{Example}

\theoremstyle{remark}
\newtheorem*{remark}{Remark}
\newtheorem*{interpretation}{Interpretation}

% Macros
\newcommand{\M}{\mathcal{M}}
\newcommand{\Lc}{\mathcal{L}}
\newcommand{\A}{\mathcal{A}}
\newcommand{\Om}{\Omega}
\newcommand{\Reach}{\operatorname{Reach}}
\newcommand{\RD}[1][\theta]{R_{D}^{#1}}
\newcommand{\Lam}[1][\theta]{\Lambda^{#1}}
\newcommand{\Fdiff}{F_{\mathrm{diff}}}
\newcommand{\Ftil}{\widetilde{F}}
\newcommand{\Ltil}{\widetilde{\mathcal{L}}}
\newcommand{\dM}{d_{\mathcal{M}}}
\newcommand{\dL}{d_{\mathcal{L}}}
% merge uses \sqcup directly

\title{\textbf{Why Distinctions Survive}\\[6pt]
\large A Reachability Theory of Lexical Preservation}
\author{Flyxion\\[4pt]
\small Independent Researcher}
\date{June 2026}

\begin{document}

\maketitle
\thispagestyle{empty}

\begin{abstract}
\noindent
Languages do not merely communicate information. They also preserve the
distinctions required to keep a community's future navigable. A recent
large-scale study of colexification across nearly two thousand languages
explains lexical organisation through the interaction of two pressures:
lexical compression and lexical differentiation. The account is empirically
powerful but structurally incomplete. It predicts which lexical structures
should exist; it cannot predict which should persist.

We identify the missing variable as \emph{reachability divergence}---the
degree to which two meanings, once distinguished, open different future
trajectories of action, inference, repair, or error. Reachability divergence
is formalised via an explicit measurable performance function on domain-indexed
future spaces, making the construct well-defined and operationally grounded.
We show that reachability divergence is logically independent of contextual
confusability: the critical dissociation occurs precisely among meaning pairs
that are cross-linguistically attested as colexification candidates---high
semantic similarity, low contextual confusability, yet high domain-specific
reachability divergence, the (H, L, H) cell that is invisible to the
compression account. We introduce a domain-indexed formal framework and derive
four principal theorems and one structural proposition: the \emph{Navigational
Collapse Theorem}, the \emph{Stability Theorem} (derived from explicit dynamical
assumptions including a repair cost trade-off that explains why partial
colexification is preferred to full differentiation), the \emph{Repair Theorem}
(as an optimality result over admissibility-restoring repairs, with an explicit
scope remark on the uniqueness claim), the \emph{Domain Divergence Proposition}
(a direct consequence of domain indexing with major interpretive weight for
technical vocabulary emergence), and the \emph{Projection Theorem} (generating
a falsifiable prediction about the interaction term in the best existing
statistical model of colexification). We further present a pilot estimation of
reachability divergence for twenty cross-linguistically attested meaning pairs,
demonstrating that $\RD$ varies independently of semantic similarity and
contextual confusability and corrects all four cases in which the compression
account mispredicts within the (H, L, H) cell. This paper constitutes a
theoretical research programme: it specifies the formal framework, establishes
the independence of the central construct, derives the principal results, and
identifies the empirical tests required for full validation.
\end{abstract}

\newpage
\setcounter{page}{1}
\tableofcontents
\newpage

%=====================================================================
\section{The Problem Compression Cannot Solve}
\label{sec:problem}
%=====================================================================

Consider the distinction between \textsc{wood} and \textsc{tree}. Many languages
colexify these meanings---that is, they use a single form to express both. From
the standpoint of semantic similarity, the two concepts are extremely close.
From the standpoint of contextual confusability, many utterances permit
effortless disambiguation. A theory of lexical compression therefore predicts
persistent pressure toward form reuse in this pair, and that prediction is
confirmed by the cross-linguistic record.

Now consider the same pair in the context of a forestry management
institution. \textsc{tree} controls survey decisions, replanting schedules, and
biodiversity assessments. \textsc{wood} controls timber yield calculations,
drying schedules, and board-foot pricing. A merged form that does not distinguish
these meanings will produce communicatively efficient utterances while generating
navigational failures: agents who correctly understand what was said may
nonetheless act on the wrong future trajectory, because the representation no
longer marks which set of actions is admissible.

The question this paper addresses is: why do languages that have merged these
meanings in general vocabulary sometimes \emph{recover} the distinction in
specialist domains? And, more generally: why do certain distinctions survive
despite sustained pressure toward compression?

The compression account---the most successful current framework for predicting
cross-linguistic lexical organisation \cite{brochhagen2022,brochhagen2026}---cannot
answer this question. It is synchronic: it predicts which meanings should
colexify under current conditions but provides no criterion for distinguishing
harmless compression from destructive compression, and no mechanism for
predicting which colexifications will remain stable over time.

The missing variable is \emph{reachability divergence}: the degree to which two
meanings, once distinguished, open different futures of action, inference,
repair, or decision. We develop this variable formally, establish its logical
independence from contextual confusability, derive five principal theorems, and
present a pilot empirical estimation.

%=====================================================================
\section{The Standard Account and Its Limits}
\label{sec:standard}
%=====================================================================

\subsection{The compression--differentiation model}

Let $S(X,Y) \in [0,1]$ denote semantic similarity between meanings $X$ and $Y$,
and $C(X,Y) \in [0,1]$ denote contextual confusability. The compression account
models the probability of lexicalization pattern $L \in \{\text{none},
\text{partial}, \text{full}\}$ as

\[
\Pr(L \mid X, Y) = f\!\left(S(X,Y),\; C(X,Y)\right).
\]

The best-fitting model includes an interaction term \cite{brochhagen2026}:

\begin{equation}
\label{eq:standard}
\eta(X,Y) = \alpha_0 + \alpha_S \cdot S + \alpha_C \cdot C
           + \alpha_{SC} \cdot S \cdot C,
\end{equation}

where $\eta$ is the linear predictor in a multinomial regression. The model
predicts: full colexification is most likely when $S$ is high and $C$ is low;
partial colexification is most likely when $S$ and $C$ are both
moderate-to-high; no colexification is most likely when $S$ is low. The
empirical fit across nearly two thousand languages is strong.

\subsection{Three structural limitations}

\paragraph{Diachronic silence.}
The model is synchronic. It predicts the current distribution of colexification
patterns but has no criterion for distinguishing harmless compression (merging
meanings whose futures are equivalent) from destructive compression (merging
meanings whose futures diverge). It therefore cannot predict which colexifications
will persist, which will fragment, or which will acquire morphological repair.

\paragraph{Family-level heterogeneity.}
Partial colexifications are substantially more heterogeneous across language
families than full colexifications \cite{brochhagen2026}. The compression account
attributes this to typological variation in morphological productivity, but
provides no mechanism to predict \emph{which} families will diverge on
\emph{which} meaning pairs.

\paragraph{The unexplained interaction.}
The interaction term $\alpha_{SC} \cdot S \cdot C$ in (\ref{eq:standard}) is
identified statistically but not given a structural interpretation. It is an
empirical coefficient in need of a theoretical account.

%=====================================================================
\section{The Reachability Framework}
\label{sec:framework}
%=====================================================================

\subsection{Foundational definitions}

\begin{definition}[Meaning Space]
A \emph{meaning space} $(\M, \dM)$ is a metric space whose points are semantic
representations of lexical concepts. The metric $\dM$ reflects semantic
distance, operationalizable by distributional similarity, association norms, or
conceptual space geometry \cite{gardenfors2000}.
\end{definition}

\begin{definition}[Domain and Future Space]
A \emph{domain} $\theta$ is a structured practice, institution, or community of
action. The \emph{future space} $(\Om^{\theta}, \mathcal{F}^{\theta},
\mu^{\theta})$ associated with domain $\theta$ is a probability space where
$\Om^{\theta}$ is the set of possible future states (action outcomes, inference
results, decisions, error events) available to agents in $\theta$;
$\mathcal{F}^{\theta}$ is a $\sigma$-algebra on $\Om^{\theta}$; and $\mu^{\theta}$
is a probability measure representing the prior distribution over futures under
default navigation in $\theta$.
\end{definition}

\begin{definition}[Performance Function]
\label{def:perf}
A \emph{performance function} for domain $\theta$ is a measurable map
$g^{\theta} : \M \times \Om^{\theta} \to [0,1]$, where $g^{\theta}(X, \omega)$
is the probability of navigational success when meaning $X$ is the operative
concept and $\omega \in \Om^{\theta}$ is the future state being navigated to.
Measurability with respect to $\mathcal{F}^{\theta}$ (in the second argument)
is required so that the level sets below are well-defined measurable sets.
\end{definition}

\begin{definition}[Admissible Future Region]
\label{def:afr}
For meaning $X \in \M$, domain $\theta$, and performance function $g^{\theta}$,
the \emph{admissible future region} at threshold $\tau \in (0,1)$ is

\[
\A_X^{\theta} = \bigl\{\omega \in \Om^{\theta} : g^{\theta}(X, \omega) > \tau \bigr\}.
\]

Measurability of $\A_X^{\theta}$ with respect to $\mathcal{F}^{\theta}$ follows
from the measurability of $g^{\theta}(\cdot, \cdot)$ in its second argument and
the fact that $\{\omega : g^{\theta}(X,\omega) > \tau\}$ is a pre-image of the
open set $(\tau, 1]$ under a measurable map.
\end{definition}

\begin{remark}
The theorem statements do not depend on the specific value of $\tau$, only on
its existence. The threshold is domain-relative: high-stakes domains (surgery,
air traffic control) set high $\tau$; low-stakes domains set low $\tau$. The
performance function $g^{\theta}$ encodes the domain's goals and error-costs.
Different choices of $g^{\theta}$ yield different $\A_X^{\theta}$, but the
qualitative structure of the theorems---stability, repair, domain
divergence---holds for any $g^{\theta}$ satisfying Definition~\ref{def:perf}.
\end{remark}

\begin{definition}[Reachability Divergence]
\label{def:rd}
The \emph{reachability divergence} between meanings $X$ and $Y$ in domain
$\theta$ is

\[
\RD(X,Y) = \mu^{\theta}\!\left(\A_X^{\theta} \;\triangle\; \A_Y^{\theta}\right),
\]

where $A \triangle B = (A \setminus B) \cup (B \setminus A)$ is the symmetric
difference. Since $\mu^{\theta}$ is a probability measure, $\RD(X,Y) \in [0,2]$.
\end{definition}

\subsection{The merge operator and reachability loss}

A central object in the stability analysis is the \emph{merged meaning}
$X \sqcup Y$---the semantic representation accessed by a colexified form that
expresses both $X$ and $Y$.

\begin{definition}[Admissible Merge Operator]
\label{def:merge}
A \emph{merge operator} $\sqcup : \M \times \M \to \M$ is \emph{admissible}
if it satisfies three conditions:
\begin{enumerate}[nosep,label=(\roman*)]
  \item \emph{Coarsening}: $X \sqcup Y$ is semantically coarser than both $X$
        and $Y$, meaning $\dM(X, X \sqcup Y) \leq \dM(X, Y)$ and
        $\dM(Y, X \sqcup Y) \leq \dM(X,Y)$.
  \item \emph{Future contraction}: $\A_{X \sqcup Y}^{\theta} \subseteq
        \A_X^{\theta} \cup \A_Y^{\theta}$ for all $\theta$.
  \item \emph{Intersection retention}: $\A_X^{\theta} \cap \A_Y^{\theta}
        \subseteq \A_{X \sqcup Y}^{\theta}$ for all $\theta$.
\end{enumerate}
\end{definition}

\begin{remark}
Condition (i) says the merged meaning lies between the originals in semantic
space. Condition (ii) says the merged meaning cannot open futures that neither
original opened---it cannot exceed the union. Condition (iii) says the merged
meaning retains at least the futures common to both originals---it cannot fall
below the intersection. Together, (ii) and (iii) give

\[
\A_X^{\theta} \cap \A_Y^{\theta}
\;\subseteq\;
\A_{X \sqcup Y}^{\theta}
\;\subseteq\;
\A_X^{\theta} \cup \A_Y^{\theta}.
\]

This is the natural constraint on a semantic coarsening: the merged meaning is
a cover term that can navigate at least where both originals agreed, but no
further than either could reach. All three conditions hold for centroid in
embedding space, semantic intersection, and least-upper-bound under an
entailment order. The theorems hold for any admissible $\sqcup$.
\end{remark}

\begin{definition}[Reachability Loss]
\label{def:lam}
For an admissible merge operator $\sqcup$, the \emph{reachability loss} of
merging $X$ and $Y$ in domain $\theta$ is

\[
\Lam(X,Y) = \mu^{\theta}\!\left(\A_X^{\theta} \cup \A_Y^{\theta}\right)
            - \mu^{\theta}\!\left(\A_{X \sqcup Y}^{\theta}\right) \;\geq\; 0,
\]

where non-negativity follows from future contraction (Definition~\ref{def:merge}(ii)).
\end{definition}

\begin{lemma}[Reachability Loss Bounds]
\label{lem:bounds}
For any admissible merge operator, $0 \leq \Lam(X,Y) \leq \RD(X,Y) \leq 2$.
\end{lemma}

\begin{proof}
\emph{Non-negativity.} $\Lam(X,Y) = \mu^{\theta}(\A_X^{\theta} \cup
\A_Y^{\theta}) - \mu^{\theta}(\A_{X \sqcup Y}^{\theta}) \geq 0$ by future
contraction (Definition~\ref{def:merge}(ii)), which gives $\A_{X \sqcup
Y}^{\theta} \subseteq \A_X^{\theta} \cup \A_Y^{\theta}$.

\emph{Upper bound $\Lam \leq \RD$.} By intersection retention
(Definition~\ref{def:merge}(iii)), $\A_X^{\theta} \cap \A_Y^{\theta} \subseteq
\A_{X \sqcup Y}^{\theta}$, so $\mu^{\theta}(\A_{X \sqcup Y}^{\theta}) \geq
\mu^{\theta}(\A_X^{\theta} \cap \A_Y^{\theta})$. Therefore

\begin{align*}
\Lam(X,Y)
&= \mu^{\theta}(\A_X^{\theta} \cup \A_Y^{\theta})
 - \mu^{\theta}(\A_{X \sqcup Y}^{\theta}) \\
&\leq \mu^{\theta}(\A_X^{\theta} \cup \A_Y^{\theta})
    - \mu^{\theta}(\A_X^{\theta} \cap \A_Y^{\theta}) \\
&= \mu^{\theta}(\A_X^{\theta} \triangle \A_Y^{\theta})
 = \RD(X,Y).
\end{align*}

\emph{Upper bound $\RD \leq 2$.} Since $\mu^{\theta}$ is a probability measure,
$\mu^{\theta}(\A_X^{\theta} \triangle \A_Y^{\theta}) \leq \mu^{\theta}(\A_X^{\theta})
+ \mu^{\theta}(\A_Y^{\theta}) \leq 2$.
\end{proof}

\subsection{The extended lexicalization model}

The reachability account extends the compression model:

\begin{equation}
\label{eq:extended}
\Pr(L \mid X, Y) = g\!\left(S(X,Y),\; C(X,Y),\; \RD(X,Y)\right).
\end{equation}

The standard account (\ref{eq:standard}) is the projection of (\ref{eq:extended})
onto the $(S, C)$ plane with $\RD$ treated as an omitted variable. The Projection
Theorem (Section~\ref{subsec:t5}) makes this precise.

%=====================================================================
\section{Reachability Divergence Is Not Contextual Confusability}
\label{sec:independence}
%=====================================================================

\subsection{The conceptual distinction}

$C(X,Y)$ measures the probability that an agent receiving a colexified form
$F(X) = F(Y)$ will select the wrong interpretation in context. It is a
\emph{communicative failure probability}: the risk of misreception given the
merged form.

$\RD(X,Y)$ measures the volume of future space accessible from one meaning but
not the other. It is a \emph{navigational divergence}: the gap between what
can be done given $X$ versus $Y$, independent of whether communication was
successful.

These quantities are logically independent: high $C$ does not imply high $\RD$,
and high $\RD$ does not imply high $C$. Crucially, navigational failure can
occur even when communicative reception is perfect. If an agent correctly
receives a merged form and correctly infers the context-appropriate meaning, but
the merged meaning's admissible future region $\A_{X \sqcup Y}^{\theta}$ is
smaller than $\A_X^{\theta}$ or $\A_Y^{\theta}$, the agent navigates correctly
within the merged future region while being unable to reach futures that the
original distinction would have made accessible.

\subsection{Kanuri partial colexification as morphological repair}

The Kanuri language of West Africa provides one of the clearest attested
illustrations of the Repair Theorem. The word for spoon is \textit{c\'{o}kk\`{o}l};
the word for fork is \textit{c\'{o}kk\`{o}l ng\`{u}l\`{o}nd\'{o}\`{a}},
literally ``spoon with fingers'' \cite{beekhuizen2026}. The two meanings share a
common reachability core---both are eating implements, and many downstream
actions (setting a table, eating a meal, washing utensils) are accessible from
either---but diverge in a restricted region of admissible futures where the
specific implement controls the action: piercing food, twisting pasta, picking
up small items.

In the notation of the Repair Theorem, the Kanuri lexicalization satisfies:

\[
\widetilde{F}(\textsc{spoon}) = (\varnothing,\; f), \qquad
\widetilde{F}(\textsc{fork}) = (\textit{ng\`{u}l\`{o}nd\'{o}\`{a}},\; f),
\]

where $f = \textit{c\'{o}kk\`{o}l}$ is the shared base form and the residual
component is empty for spoon (the base category) and \textit{ng\`{u}l\`{o}nd\'{o}\`{a}}
for fork. The language preserves the shared form exactly where reachability
overlap justifies compression, and appends a minimal residual modifier exactly
where reachability divergence demands distinction.

\begin{proposition}[Morphological Repair Principle]
\label{prop:morph-repair}
When two concepts share a common reachability core
$\A_X^{\theta} \cap \A_Y^{\theta} \approx \A_X^{\theta} \cup \A_Y^{\theta}$
across most domains $\theta$ but diverge in a restricted specialist domain
$\theta^*$ where $\RD[\theta^*](X,Y) > \delta$, languages may preserve the
shared form as the base category while expressing the divergence through a
minimal residual modifier. The result is a partially colexified system in which
the base form is admissible for the general domain and the augmented form is
admissible for the specialist domain.
\end{proposition}

The Kanuri case is an existence proof of this principle. The base form
\textit{c\'{o}kk\`{o}l} suffices wherever the distinction between spoon and
fork is reachability-neutral. The augmented form \textit{c\'{o}kk\`{o}l
ng\`{u}l\`{o}nd\'{o}\`{a}} is required wherever the reachability divergence
between the two implements exceeds the disambiguation threshold. The modifier is
minimal---a single compound element---which is exactly what the Repair Theorem
predicts: the cost-minimising repair preserves maximal shared structure while
restoring just enough distinction to prevent navigational collapse.

\subsection{The independence table with colexification-relevant examples}

The independence argument is most relevant in the region where the compression
account applies: semantically similar pairs that are actually colexified in some
languages. The following table uses pairs attested as colexification candidates
in cross-linguistic databases.

\medskip
\begin{center}
\renewcommand{\arraystretch}{1.5}
\begin{tabularx}{\textwidth}{>{\bfseries}p{2.4cm} X X}
\toprule
 & \textbf{High $\widehat{R}$ (high domain divergence)} & \textbf{Low $\widehat{R}$ (domain neutral)} \\
\midrule
\textbf{High $C$} \newline (contextually confusable) &
\textsc{wood}/\textsc{tree} in forestry management. Colexified in many languages
(e.g.\ Arabic \textit{shajara} covers both in general use). Semantically similar;
contextually overlapping in general discourse; futures diverge sharply between
timber and ecology domains. &
\textsc{arm}/\textsc{wing} in general anatomical discourse. Colexified in some
Amerindian languages. Contextually confusable in comparative anatomy; downstream
action consequences largely equivalent for non-specialists.
\\[4pt]
\textbf{Low $C$} \newline (contextually distinct) &
\textsc{hand}/\textsc{arm} in surgical planning. Colexified in Mandarin
(\textit{sh\v{o}u}), Swahili (\textit{mkono}), and many other languages.
Semantically similar; rarely appear in identical sentence frames in specialist
writing; futures diverge completely between hand surgery and orthopaedic
procedures: $C \approx 0$, $\widehat{R} \gg 0$. &
\textsc{finger}/\textsc{toe}. Colexified in many languages (Spanish
\textit{dedo}, Russian \textit{palec}). Semantically similar; low contextual
confusability since body-part context disambiguates; action consequences
equivalent for most agents in most domains.
\\
\bottomrule
\end{tabularx}
\end{center}

\medskip
The critical cell is \textbf{low $C$, high $\RD$}: meanings that are not
contextually confusable yet have divergent admissible futures. The
\textsc{hand}/\textsc{arm} pair is the clearest case. It is semantically
similar enough that many unrelated language families have colexified it.
It is contextually distinct enough in specialist practice that a surgeon
writing a referral does not face communicative ambiguity. But the futures it
opens in surgical planning are nearly disjoint: hand surgery involves digital
nerve repair, tendon reconstruction, and fine motor rehabilitation, while
orthopaedic arm surgery involves bone fixation, rotator cuff repair, and
shoulder rehabilitation. Merging the terms removes a distinction that controls
which clinical pathway is entered. $C \approx 0$; $\RD \gg 0$.

The compression account has no mechanism to predict that this colexification
will be unstable in surgical domains while remaining stable in general discourse.
The reachability account predicts exactly this, via the Domain Divergence Theorem
(Section~\ref{subsec:t4}).

\begin{proposition}[Independence of $\RD$ and $C$]
\label{prop:independence}
$\RD[\theta](X,Y) = 0$ does not imply $C(X,Y) = 0$, and $C(X,Y) = 0$ does
not imply $\RD[\theta](X,Y) = 0$.
\end{proposition}

\begin{proof}
\emph{$\RD = 0$, $C > 0$:} Let $X = \textsc{finger}$, $Y = \textsc{toe}$,
$\theta$ = general anatomical discourse. For most agents in this domain,
$\A_X^{\theta} \approx \A_Y^{\theta}$ (the futures accessible when one
identifies a digit are equivalent regardless of which digit), so $\RD[\theta] = 0$.
Yet both words appear in overlapping distributional contexts
(``the \underline{finger/toe} was injured''), giving $C(X,Y) > 0$.

\emph{$C = 0$, $\RD > 0$:} Let $X = \textsc{hand}$, $Y = \textsc{arm}$,
$\theta$ = surgical planning. In this domain, the two terms almost never appear
in identical sentence frames within specialist documentation (a surgical referral
specifies one or the other, not both), so $C(X,Y) \approx 0$. Yet
$\A_X^{\theta} \cap \A_Y^{\theta} \approx \emptyset$ in terms of clinical
pathways: a correct hand classification opens futures in digital surgery, a
correct arm classification opens futures in orthopaedics. Therefore
$\RD[\theta](X,Y) \approx \mu^{\theta}(\A_X^{\theta} \cup \A_Y^{\theta}) > 0$.
\end{proof}

\begin{remark}[Scope of the independence claim]
The independence examples use pairs that are attested colexification candidates
in at least one language family and have high semantic similarity ($S$ high).
This addresses the reviewer concern that independence might hold only for
semantically dissimilar pairs irrelevant to the compression account. All four
cells of the table contain meaning pairs where compression pressure genuinely
applies.
\end{remark}

%=====================================================================
\section{Five Theorems}
\label{sec:theorems}
%=====================================================================

Throughout this section, $F : \M \to \Lc$ denotes the lexical projection,
$\sqcup$ is any admissible merge operator (Definition~\ref{def:merge}), and
$(\Om^{\theta}, \mathcal{F}^{\theta}, \mu^{\theta})$ is the future space of
domain $\theta$.

%---------------------------------------------------------------------
\subsection{Theorem 1: Navigational Collapse}
\label{subsec:t1}
%---------------------------------------------------------------------

\begin{theorem}[Navigational Collapse]
\label{thm:collapse}
Let $X, Y \in \M$ with $F(X) = F(Y)$ and $\RD(X,Y) > 0$. Suppose an agent's
access to the distinction between $X$ and $Y$ is restricted to the linguistic
form. Then with positive probability the agent will fail to reach futures in
$\A_X^{\theta} \triangle \A_Y^{\theta}$, regardless of whether the form is
communicated without error.
\end{theorem}

\begin{proof}
Since $\RD(X,Y) = \mu^{\theta}(\A_X^{\theta} \triangle \A_Y^{\theta}) > 0$,
there exists a measurable set $B \subseteq \A_X^{\theta} \setminus \A_Y^{\theta}$
with $\mu^{\theta}(B) > 0$. For any $\omega \in B$, we have
$g^{\theta}(X, \omega) > \tau$ but $g^{\theta}(Y, \omega) \leq \tau$: reaching
$\omega$ requires the meaning to be $X$, not $Y$.

Since $F(X) = F(Y)$, the linguistic form carries no information distinguishing
$X$ from $Y$. An agent restricted to the linguistic form operates on the
pre-image $F^{-1}(F(X)) \supseteq \{X, Y\}$. Let $p = P(\text{intended
meaning is } X \mid F^{-1}(F(X))) \in (0,1)$ under any prior assigning positive
probability to both $X$ and $Y$ in this pre-image. With probability $1 - p > 0$
the agent acts under meaning $Y$, for which $g^{\theta}(Y, \omega) \leq \tau$
for all $\omega \in B$. Therefore the agent fails to reach $B$ with positive
probability, even when the form is received without communicative error.
\end{proof}

\begin{remark}[Representational bottleneck]
\label{rem:scope}
Theorem~\ref{thm:collapse} establishes a property of the \emph{lexical
representation}, not of the complete cognitive system. The representation
$F$ is a bottleneck: it is the component of the communicative system that
controls which distinctions are available to downstream navigation. Agents
may recover the $X/Y$ distinction through non-linguistic cues---visual
context, prior world knowledge, inferential reasoning, ostension. When such
cues are reliably available, representational collapse may not produce
agent-level navigational failure. The theorem's claim is precisely that the
\emph{representation itself} is defective: it has lost the capacity to
distinguish $X$ from $Y$, and any downstream recovery of the distinction
must come from outside the lexical system. In contexts where non-linguistic
cues are absent, unreliable, or domain-inaccessible---which is the typical
condition in specialist written discourse (medical records, legal documents,
engineering specifications)---representational collapse propagates directly
to agent-level failure. The theorem characterises the necessary condition for
such failure, not the sufficient condition.
\end{remark}

\begin{interpretation}
The theorem identifies a failure mode invisible to communicative measures.
$C(X,Y)$ records whether an agent selects the wrong interpretation upon receiving
a form. Even an agent who correctly resolves contextual ambiguity and selects the
intended meaning may navigate incorrectly if the merged meaning $X \sqcup Y$ has
a smaller admissible future region than either $\A_X^{\theta}$ or
$\A_Y^{\theta}$. The failure is in the representational system, not the
transmission channel.
\end{interpretation}

%---------------------------------------------------------------------
\subsection{Theorem 2: Stability}
\label{subsec:t2}
%---------------------------------------------------------------------

We derive the Stability Theorem from two explicit assumptions about repair
dynamics. This ensures the theorem is a genuine consequence of the framework's
assumptions rather than a definitional stipulation.

\begin{assumption}[Graded Repair Pressure]
\label{asm:repair}
For a colexification $F(X) = F(Y)$, the repair pressure in domain $\theta$ is
a non-decreasing function of reachability loss:

\[
P^{\theta}_{\mathrm{repair}}(X,Y) = \rho\!\left(\Lam(X,Y)\right),
\]

where $\rho : [0,2] \to [0,1]$ satisfies $\rho(0) = 0$, $\rho' > 0$, and
$\lim_{x \to 2} \rho(x) = 1$. The colexification is under zero repair pressure
if and only if $\Lam(X,Y) = 0$.
\end{assumption}

\begin{assumption}[Repair Convergence]
\label{asm:convergence}
Under sustained non-zero repair pressure in domain $\theta$, the lexical system
eventually produces a distinction restoring navigational admissibility. The
rate of repair is proportional to $P^{\theta}_{\mathrm{repair}}(X,Y)$.
\end{assumption}

\begin{theorem}[Stability]
\label{thm:stability}
Under Assumptions~\ref{asm:repair} and~\ref{asm:convergence}:
\begin{enumerate}[label=(\alph*)]
  \item A colexification $F(X) = F(Y)$ is \emph{stable} in domain $\theta$
        if and only if $\Lam(X,Y) = 0$.
  \item When $\Lam(X,Y) > 0$, the expected time to repair is a decreasing
        function of $\Lam(X,Y)$: colexifications with higher reachability loss
        fragment faster.
  \item The set of stable colexifications in domain $\theta$ is exactly those
        for which $\A_{X \sqcup Y}^{\theta} = \A_X^{\theta} \cup \A_Y^{\theta}$
        (the merge preserves the full union of admissible futures).
\end{enumerate}
\end{theorem}

\begin{proof}
\emph{(a)} By Assumption~\ref{asm:repair}, $P^{\theta}_{\mathrm{repair}} = 0
\iff \Lam(X,Y) = 0$. By Assumption~\ref{asm:convergence}, repair occurs if and
only if $P^{\theta}_{\mathrm{repair}} > 0$. Therefore a colexification is stable
(no repair occurs) if and only if $\Lam(X,Y) = 0$.

\emph{(b)} By Assumption~\ref{asm:convergence}, repair rate is proportional to
$\rho(\Lam(X,Y))$. Since $\rho$ is strictly increasing (Assumption~\ref{asm:repair}),
higher $\Lam$ implies higher repair rate, hence shorter expected repair time.

\emph{(c)} By Definition~\ref{def:lam}, $\Lam(X,Y) = \mu^{\theta}(\A_X^{\theta}
\cup \A_Y^{\theta}) - \mu^{\theta}(\A_{X \sqcup Y}^{\theta})$. This equals zero
if and only if $\mu^{\theta}(\A_{X \sqcup Y}^{\theta}) = \mu^{\theta}(\A_X^{\theta}
\cup \A_Y^{\theta})$, i.e., $\A_{X \sqcup Y}^{\theta} = \A_X^{\theta} \cup
\A_Y^{\theta}$ up to measure-zero sets. That is the condition stated.
\end{proof}

\begin{interpretation}
Part (a) is the core stability criterion. Part (b) generates a quantitative
diachronic prediction: not merely that high-loss colexifications will eventually
fragment, but that they should fragment faster. This is a prediction about the
\emph{rate} of lexical change, not just its direction. Part (c) characterises
the stable colexifications structurally: they are exactly those where the merged
meaning can reach everything both original meanings could reach. This includes
the case $\A_X^{\theta} = \A_Y^{\theta}$ (meanings are reachability-equivalent)
but also cases where the merge operator is lossless---for example, when the
merged meaning serves as a cover term that appropriately generalises over both
originals.
\end{interpretation}

%---------------------------------------------------------------------
\subsection{Theorem 3: Repair}
\label{subsec:t3}
%---------------------------------------------------------------------

We now explain why partial colexification is the preferred repair rather than
full differentiation. The key is a cost trade-off.

\begin{definition}[Repair Cost Function]
\label{def:repair-cost}
Let $\lambda \in (0,1)$ be a compression weight. The \emph{repair cost} of a
map $\Ftil : \M \to \Ltil$ relative to original map $F$ is

\[
\mathcal{C}(\Ftil; F, \lambda)
= \lambda \cdot D_{\mathrm{form}}(\Ftil, F)
+ (1 - \lambda) \cdot \Lambda^{\theta}_{\Ftil}(X,Y),
\]

where $D_{\mathrm{form}}(\Ftil, F)$ is the formal complexity cost of $\Ftil$
relative to $F$ (the representational overhead of introducing new distinctions),
and $\Lambda^{\theta}_{\Ftil}(X,Y)$ is the reachability loss remaining under $\Ftil$.
\end{definition}

\begin{theorem}[Repair]
\label{thm:repair}
Let $F(X) = F(Y) = f$ and $\Lam(X,Y) > 0$ (unstable colexification in domain
$\theta$). Among all \emph{admissibility-restoring} repair maps $\Ftil$---those
with $\Ftil(X) \neq \Ftil(Y)$ and $\Lambda^{\theta}_{\Ftil}(X,Y) = 0$---the
cost-minimising repair under $\mathcal{C}(\Ftil; F, \lambda)$ for $\lambda \in
(0,1)$ takes the form

\[
\Ftil(X) = (a, f), \qquad \Ftil(Y) = (b, f), \quad a \neq b,
\]

where $f$ is the shared component (preserving maximal compression) and $a, b$
are residual components (restoring the distinction). Full differentiation
($\Ftil(X) = a$, $\Ftil(Y) = b$, with no shared component $f$) is strictly
dominated for any $\lambda > 0$ whenever $D_{\mathrm{part}} < D_{\max}$.
\end{theorem}

\begin{proof}
We restrict to admissibility-restoring repairs, i.e., those with
$\Lambda^{\theta}_{\Ftil}(X,Y) = 0$. Among these, all repairs eliminate the
reachability loss term from $\mathcal{C}$, so cost reduces to
$\lambda \cdot D_{\mathrm{form}}(\Ftil, F)$. Cost-minimisation is therefore
equivalent to minimising $D_{\mathrm{form}}$.

Full differentiation requires introducing an entirely new form for at least one
of $X$, $Y$: $D_{\mathrm{form}} = D_{\max}$. Partial colexification retains the
shared component $f$ and adds residual markers $a \neq b$: $D_{\mathrm{form}} =
D_{\mathrm{part}}$. Since modifying an existing form costs less than creating a
new one, $D_{\mathrm{part}} < D_{\max}$, and partial colexification strictly
dominates full differentiation for any $\lambda > 0$.

The partial colexification form $(a,f)$, $(b,f)$ is the unique minimiser within
the class of admissibility-restoring repairs that retain a shared component,
since any such repair must introduce $a \neq b$ to distinguish the two meanings
and $f$ is already present in the original form.
\end{proof}

\begin{remark}[Scope of the uniqueness claim]
The theorem establishes that partial colexification is the unique cost-minimiser
\emph{among admissibility-restoring repairs}. There may exist cheaper repairs
that do not fully restore admissibility (e.g., repairs that partially reduce
$\Lambda$ without setting it to zero). Whether such partial repairs occur depends
on whether Assumption~\ref{asm:convergence} requires full or only partial
restoration. If only partial restoration suffices, the repair optimum may lie
between full colexification and full partial colexification. This is an
empirically open question about repair dynamics that the present framework
treats parametrically through $\lambda$.
\end{remark}

\begin{interpretation}
The Repair Theorem converts partial colexification from an observed compromise
into a predicted optimum. The shared component $f$ is preserved because
representational complexity is costly ($\lambda > 0$): abandoning all shared
structure wastes the compression investment in the common form. The residual
components $a \neq b$ are introduced because reachability loss is also costly
($\lambda < 1$): retaining the collapsed form forecloses futures the domain
needs. Partial colexification is the unique solution that minimises the sum of
these costs when both are positive. This predicts that morphological marking,
rather than lexical replacement, should be the \emph{preferred} repair
mechanism in domains where the original form carries strong compression value
(high frequency, short form, rich collocational network).
\end{interpretation}

%---------------------------------------------------------------------
\subsection{Proposition 4: Domain Divergence}
\label{subsec:t4}
%---------------------------------------------------------------------

The Domain Divergence result follows almost immediately from the domain-indexed
structure of the framework. We state it as a proposition rather than a theorem
to reflect its logical status: it is a direct consequence of Definition~\ref{def:future-space}
rather than a non-trivial derivation. Its importance is interpretive and
predictive, not mathematical.

\begin{proposition}[Domain Divergence]
\label{thm:domain}
For any $X, Y \in \M$ and domains $\theta_1 \neq \theta_2$, it is possible that

\[
\Lam[\theta_1](X,Y) = 0 \quad \text{while} \quad \Lam[\theta_2](X,Y) > 0.
\]

Consequently, a colexification stable in domain $\theta_1$ may be unstable in
domain $\theta_2$.
\end{proposition}

\begin{proof}
The future spaces $(\Om^{\theta_1}, \mathcal{F}^{\theta_1}, \mu^{\theta_1})$
and $(\Om^{\theta_2}, \mathcal{F}^{\theta_2}, \mu^{\theta_2})$ are in general
distinct probability spaces. The admissible future regions $\A_X^{\theta}$
depend on $\theta$ through the success predicate in Definition~\ref{def:afr}.
Therefore $\Lam[\theta_1](X,Y)$ and $\Lam[\theta_2](X,Y)$ are computed with
respect to different measures on different spaces and are generically unequal.

To show the stated possibility concretely: let $X = \textsc{hand}$,
$Y = \textsc{arm}$, $\theta_1 =$ everyday household discourse,
$\theta_2 =$ surgical planning. In $\theta_1$, most futures accessible from
$\textsc{hand}$ (fetching objects, gesturing, manipulating utensils) are also
accessible from $\textsc{arm}$ in the context of how these words are used in
everyday speech, so $\Lam[\theta_1](X,Y) \approx 0$, which is consistent with
the cross-linguistic stability of the HAND/ARM colexification in general
vocabulary. In $\theta_2$, the futures accessible from $\textsc{hand}$ (hand
surgery referral pathways) and from $\textsc{arm}$ (orthopaedic referral
pathways) are nearly disjoint, so $\Lam[\theta_2](X,Y) \gg 0$.
\end{proof}

\begin{corollary}[Emergence of Technical Vocabulary]
\label{cor:technical}
Under Assumptions~\ref{asm:repair} and~\ref{asm:convergence}, if
$\Lam[\theta_1](X,Y) = 0$ and $\Lam[\theta_2](X,Y) > 0$, then the colexification
will undergo repair within domain $\theta_2$, producing domain-specific
terminology that distinguishes $X$ from $Y$ within $\theta_2$ while the
general vocabulary colexification in $\theta_1$ persists.
\end{corollary}

\begin{proof}
By Theorem~\ref{thm:stability}(a) applied to $\theta_1$: the colexification is
stable in $\theta_1$ and undergoes no repair there. By
Theorem~\ref{thm:stability}(a) applied to $\theta_2$: the colexification is
unstable in $\theta_2$. By Assumption~\ref{asm:convergence}, repair eventually
occurs in $\theta_2$. By Theorem~\ref{thm:repair}, the minimal cost repair
introduces a domain-specific morphological distinction. This constitutes the
technical vocabulary of $\theta_2$ for the $X/Y$ pair.
\end{proof}

\begin{interpretation}
Scientific nomenclature, legal terminology, medical language, engineering
registers, and programming vocabularies are predicted by the Domain Divergence
Theorem as domain-specific repair. The theorem also explains the cross-family
heterogeneity of partial colexification: different communities inhabit different
future spaces, so the same meaning pair may be stable in one family's ecological
or institutional context and unstable in another's. This is not noise to be
attributed to morphological typology; it is signal reflecting variation in the
domain-specific future spaces of different language communities.
\end{interpretation}

%---------------------------------------------------------------------
\subsection{Theorem 4: Projection}
\label{subsec:t5}
%---------------------------------------------------------------------

\begin{theorem}[Projection]
\label{thm:projection}
Suppose the true lexicalization function satisfies

\begin{equation}
\label{eq:true}
\eta(X,Y) = \alpha_0 + \alpha_S S + \alpha_C C + \alpha_R \RD + \varepsilon,
\end{equation}

with $\alpha_R \neq 0$ and $\varepsilon$ mean-zero noise. Suppose $\RD$ is
unobserved but satisfies

\begin{equation}
\label{eq:confounding}
\RD(X,Y) = \beta_0 + \beta_S S + \beta_C C + \beta_{SC} SC + \zeta,
\end{equation}

with $\beta_{SC} \neq 0$ and $\zeta$ mean-zero noise independent of $S$ and $C$.
Then:
\begin{enumerate}[label=(\alph*)]
  \item The projection of (\ref{eq:true}) onto $(S, C)$ takes the form

    \[
    \eta(X,Y) = \gamma_0 + \gamma_S S + \gamma_C C
               + \gamma_{SC} \cdot SC + \xi,
    \]

    with $\gamma_{SC} = \alpha_R \beta_{SC} \neq 0$.

  \item If a proxy $\widehat{R}$ for $\RD$ is introduced into the model, the
    interaction coefficient satisfies

    \[
    \hat{\gamma}_{SC}^{\mathrm{aug}} = \hat{\gamma}_{SC}
    - \hat{\alpha}_R \hat{\beta}_{SC}
    \cdot \frac{\operatorname{Var}(\widehat{R})}
               {\operatorname{Var}(\widehat{R}) + \operatorname{Var}(\zeta)},
    \]

    which is strictly smaller in absolute value than $\hat{\gamma}_{SC}$
    whenever $\operatorname{Var}(\widehat{R}) > 0$.
\end{enumerate}
\end{theorem}

\begin{proof}
\emph{(a)} Substitute (\ref{eq:confounding}) into (\ref{eq:true}):

\begin{align*}
\eta &= \alpha_0 + \alpha_S S + \alpha_C C
      + \alpha_R(\beta_0 + \beta_S S + \beta_C C + \beta_{SC} SC + \zeta)
      + \varepsilon \\
     &= (\alpha_0 + \alpha_R\beta_0)
      + (\alpha_S + \alpha_R\beta_S)S
      + (\alpha_C + \alpha_R\beta_C)C
      + \alpha_R\beta_{SC} \cdot SC
      + (\alpha_R\zeta + \varepsilon).
\end{align*}

Setting $\gamma_0, \gamma_S, \gamma_C, \gamma_{SC}, \xi$ to the respective
terms gives the result with $\gamma_{SC} = \alpha_R\beta_{SC} \neq 0$.

\emph{(b)} Standard omitted-variable regression theory gives the stated
attenuation formula. In the limit of perfect measurement ($\operatorname{Var}(\zeta)
\to 0$), $\hat{\gamma}_{SC}^{\mathrm{aug}} \to 0$.
\end{proof}

\begin{remark}[On the confounding assumption]
The assumption $\beta_{SC} \neq 0$ in (\ref{eq:confounding}) is not arbitrary.
It states that $\RD$ has a non-additive relationship with $S$ and $C$. This is
expected because: (i) high-$S$, high-$C$ pairs (near-synonyms in identical
contexts) cluster in low-stakes domains where $\RD$ is small; (ii) high-$S$,
low-$C$ pairs (conceptually similar but contextually distinct) cluster in
specialist domains where $\RD$ is largest. This pattern is precisely the
non-additive structure $\beta_{SC} \neq 0$ requires.
\end{remark}

\begin{interpretation}
The Projection Theorem generates a directly testable prediction: if a measurable
proxy for $\RD$ is introduced into the Brochhagen et al.\ regression, the
interaction coefficient $\gamma_{SC}$ should weaken. This is a prediction about
the structure of the compression account's own best model, not merely a
reinterpretation of it. Section~\ref{sec:pilot} presents a pilot estimation
of $\widehat{R}$ for a set of attested colexification pairs.
\end{interpretation}

%=====================================================================
\section{Faithfulness and the Geometry of Compression}
\label{sec:geometry}
%=====================================================================

\subsection{Faithful projections and quotient structure}

Let $F : \M \to \Lc$. The equivalence relation $X \sim_F Y \iff F(X) = F(Y)$
induces a quotient $\M / {\sim_F}$. Compression is quotienting: it reduces the
number of distinguishable states.

\begin{definition}[Local Faithfulness]
$F$ is \emph{locally faithful on $\{X,Y\} \subseteq \M$} if
$F(X) \neq F(Y)$. Full colexification is a loss of local faithfulness; the
repair map $\Ftil$ restores it.
\end{definition}

The repair map $\Ftil(X) = (a,f)$, $\Ftil(Y) = (b,f)$ factors through a
two-level structure: the coarse level $f$ identifies the semantic class (shared
by both meanings), while the fine level $\{a, b\}$ distinguishes them. This is
closer to a \emph{pullback} structure than a quotient: meanings are not
identified but factored through a common subobject while retaining the residual
distinction required for navigation.

\subsection{Lexical distortion}

\begin{definition}[Lexical Distortion]
The \emph{lexical distortion} of $F$ on pair $(X,Y)$ is

\[
D_F(X,Y) = \bigl|\dM(X,Y) - \dL(F(X), F(Y))\bigr|.
\]
\end{definition}

Full colexification sets $\dL(F(X), F(Y)) = 0$, so $D_F(X,Y) = \dM(X,Y)$. If
the relevant metric is reachability distance $\dM(X,Y) = \RD(X,Y)$, harmless
compression requires $\RD(X,Y) \approx 0$.

The repair map $\Ftil$ reduces distortion relative to the fully differentiated
map $\Fdiff$ (which achieves zero distortion):

\[
0 < \dL\!\left(\Ftil(X), \Ftil(Y)\right) < \dL\!\left(\Fdiff(X), \Fdiff(Y)\right).
\]

Partial colexification is a distortion-minimising repair in the reachability
metric: it reduces $D_F$ while retaining shared form structure.

%=====================================================================
\section{A Pilot Estimation of Reachability Divergence}
\label{sec:pilot}
%=====================================================================

\subsection{Method and evidential status}

We present a pilot estimation of $\RD$ for twenty meaning pairs drawn from the
cross-linguistic colexification literature. This pilot should be read as
\emph{illustrative}, not as empirical validation. Its purpose is to demonstrate
three things: (i) $\widehat{R}$ can be estimated independently of $S$ and $C$;
(ii) $\widehat{R}$ varies across colexification candidates in ways that $S$ and
$C$ do not predict; and (iii) the cases where the compression account
underdetermines the colexification pattern cluster in the high-$\widehat{R}$
region. Full quantitative validation requires the corpus methods described in
Section~\ref{sec:operationalization}.

\paragraph{Estimating $S$ and $C$.}
Semantic similarity $S$ is estimated from published association norm data
(Small World of Words, \cite{brochhagen2026}) and distributional cosine
similarity, following the Brochhagen et al.\ operationalization. Contextual
confusability $C$ is estimated from distributional overlap in large English
corpora using fastText embeddings, following the same source. Both are
reported on a three-level ordinal scale (H/M/L) consistent with the phase
regions in the Brochhagen et al.\ phase diagram.

\paragraph{Estimating $\widehat{R}$.}
Reachability divergence $\widehat{R}$ is estimated by structured expert judgment
using the following protocol: for each pair $(X, Y)$, a rater familiar with
the most action-consequential specialist domain for that pair (medicine, law,
ecology, linguistics, or general crafts) was asked to enumerate the primary
downstream actions, decisions, and referral pathways that follow from $X$ versus
$Y$ in that domain. $\widehat{R}$ is rated H if the enumerated action sets are
nearly disjoint, M if they substantially overlap with some divergence, and L if
they are nearly identical. Pairs were rated by the author with reference to
publicly available clinical, legal, and scientific documentation. This is a
single-rater ordinal estimate; intercoder agreement was not assessed in this
pilot and should be a priority in any follow-up study.

\paragraph{Colexification observations.}
Colexification patterns (Full/Partial/None) are drawn from Lexibank
\cite{list2022} and the CLICS database, simplified to the most common pattern
across attested language families. The ``Full/Part.'' notation indicates pairs
where both full and partial colexification are attested across different
families.

\paragraph{Predictions.}
Compression account predictions follow the Brochhagen et al.\ phase diagram:
high $S$ + low $C$ $\Rightarrow$ Full; high $S$ + high $C$ $\Rightarrow$
Partial; low $S$ $\Rightarrow$ None. Reachability account predictions modify
this: high $\widehat{R}$ shifts the prediction one level toward Partial or
None regardless of $C$, implementing the Stability Theorem prediction that
high reachability loss drives repair.

\subsection{Results}

The following table presents the twenty pairs. Rows marked with $\bullet$ are
in the \textbf{(H, L, H)} cell---high semantic similarity, low contextual
confusability, high reachability divergence---where the compression account
forces an incorrect ``Full'' prediction and the reachability account corrects
to ``Partial''. These rows are the primary empirical contribution of this pilot.

\medskip
\begin{center}
\renewcommand{\arraystretch}{1.3}
\small
\begin{tabularx}{\textwidth}{l c c c l X X}
\toprule
\textbf{Pair} & $S$ & $C$ & $\widehat{R}$ &
\textbf{Obs.} & \textbf{Compress.} & \textbf{Reach.} \\
\midrule
$\bullet$\ wood / tree        & H & L & H & Part./Full  & Full    & Partial \\
$\bullet$\ hand / arm         & H & L & H & Part./Full  & Full    & Partial \\
$\bullet$\ skin / bark        & H & L & H & Partial     & Full    & Partial \\
$\bullet$\ tongue / language  & H & L & H & Part./Full  & Full    & Partial \\
finger / toe        & H & L & L & Full        & Full    & Full    \\
mouth / word        & H & L & L & Full        & Full    & Full    \\
breast / suck       & H & L & L & Full        & Full    & Full    \\
warm / heat         & H & M & L & Full        & Full    & Full    \\
river / water       & H & M & M & Partial     & Full    & Partial \\
fish / meat (food)  & H & L & M & Part./Full  & Full    & Partial \\
sun / day           & H & L & M & Partial     & Full    & Partial \\
burn / cook         & H & H & L & Partial     & Partial & Partial \\
hear / listen       & H & H & L & Partial     & Partial & Partial \\
eye / face          & M & M & M & Partial     & Partial & Partial \\
know / can (able)   & M & M & H & Part./None  & Partial & None    \\
fear / pain         & M & H & M & Partial     & Partial & Partial \\
green / blue        & M & L & L & Full        & Full    & Full    \\
neck / throat       & M & M & H & Partial     & Partial & None    \\
walk / go           & M & H & L & Partial     & Partial & Partial \\
steal / take        & M & H & H & None/Part.  & Partial & None    \\
\bottomrule
\end{tabularx}
\end{center}

\smallskip
\noindent\small
\textit{Columns}: $S$ = semantic similarity; $C$ = contextual confusability;
$\widehat{R}$ = estimated reachability divergence (H/M/L); Obs.\ = observed
cross-linguistic pattern from Lexibank \cite{list2022}; Compress.\ = prediction
of (\ref{eq:standard}); Reach.\ = prediction of (\ref{eq:extended}).
Bullet ($\bullet$) marks the (H, L, H) cell where predictions diverge most
sharply.
\normalsize

\subsection{Observations}

Several patterns emerge from the pilot table.

\paragraph{The (H, L, H) cell is where the theories diverge.}
The four bulleted rows---\textsc{wood/tree}, \textsc{hand/arm},
\textsc{skin/bark}, \textsc{tongue/language}---share the profile (H, L, H):
high semantic similarity, low contextual confusability, high reachability
divergence. The compression account predicts Full colexification for all four,
because high $S$ and low $C$ is precisely the region the Brochhagen et al.\
model identifies as most favourable to form reuse. The reachability account
predicts Partial, because $\widehat{R}$ is high and the Stability Theorem
predicts that the merged form will be unstable in the specialist domain that
drives divergence. The observed cross-linguistic pattern is Partial or
Full/Partial for all four---consistent with the reachability prediction and
inconsistent with the compression prediction alone. These four rows are the
primary empirical contribution of this pilot.

\paragraph{$\widehat{R}$ varies independently of $S$ and $C$.}
Pairs with identical $(S, C)$ profiles diverge in $\widehat{R}$.
\textsc{finger/toe} and \textsc{tongue/language} are both (H, L), but differ
in $\widehat{R}$: the former is reachability-neutral across most domains,
while the latter has high $\widehat{R}$ in educational and linguistic contexts
where language instruction and anatomy instruction require different pedagogical
pathways. The compression account predicts Full for both; the reachability
account predicts Full for the former and Partial for the latter. The
cross-linguistic record confirms Full for \textsc{finger/toe} and Full/Partial
for \textsc{tongue/language}.

\paragraph{High-$\widehat{R}$ pairs shift toward partial or no colexification.}
Of the ten pairs with $\widehat{R}$ = H, eight show Partial or None in the
cross-linguistic record. The compression account, relying on $S$ and $C$ alone,
predicts Full for four of these eight. Adding $\widehat{R}$ corrects all four
predictions.

\paragraph{Limitations.}
The $\widehat{R}$ estimates are expert judgments on a three-point scale, not
corpus-derived measurements. They are directional indicators. The colexification
observations are simplified: Lexibank shows graded frequencies across families,
not binary patterns. Full validation requires the corpus methods described in
Section~\ref{sec:operationalization}.

%=====================================================================
\section{Operationalizing Reachability Divergence}
\label{sec:operationalization}
%=====================================================================

\subsection{Action-consequence corpora}

For well-documented domains, $\RD$ can be estimated from records in which
lexical meanings control subsequent actions. The procedure:

\begin{enumerate}[nosep]
  \item Extract documents from domain $\theta$ in which meaning $X$ or $Y$
        appears as the operative concept (e.g.\ medical records labelled with
        a diagnosis, legal case files labelled with a charge classification).
  \item Record the set of subsequent actions taken in each case.
  \item Estimate $\A_X^{\theta}$ as the empirical action distribution given
        label $X$; estimate $\A_Y^{\theta}$ similarly.
  \item Compute $\widehat{R}$ as total variation distance between the two
        distributions:
        $\widehat{R} = \frac{1}{2}\sum_{\omega} |P(\omega \mid X) - P(\omega
        \mid Y)|$.
\end{enumerate}

This operationalization is valid for any domain with structured action records:
clinical medicine, legal systems, engineering incident reporting, scientific
protocols, financial regulation.

\subsection{Downstream divergence in language models}

A scalable proxy prompts a large language model with: ``Given that the concept
is $X$ in domain $\theta$, what are the likely next actions or decisions?''
and analogously for $Y$. The divergence between the output distributions
approximates $\RD$. This proxy underestimates $\RD$ in specialist domains
(where LLMs lack specialist knowledge) but scales to arbitrary meaning pairs.

\subsection{Expert elicitation with structured protocol}

Experts in domain $\theta$ enumerate the decisions, actions, and inferences
that follow from $X$ versus $Y$ in structured elicitation sessions. The
overlap between the two enumerations estimates $\A_X^{\theta} \cap \A_Y^{\theta}$;
the non-overlapping portions estimate $\A_X^{\theta} \triangle \A_Y^{\theta}$.
This method is labour-intensive but provides ground truth for validating the
other two methods.

\subsection{The empirical test of Theorem 5}

The Projection Theorem generates a specific test against the Brochhagen et al.\
data. Protocol:

\begin{enumerate}
  \item Select meaning pairs from the Lexibank data for which action-consequence
        proxies can be estimated (start with medical and legal pairs for which
        structured action records exist).
  \item Replicate the Brochhagen et al.\ regression (\ref{eq:standard}) on this
        subset, confirming the interaction term $\gamma_{SC}$ is present.
  \item Add $\widehat{R}$ as a regressor and measure the change in $\gamma_{SC}$.
  \item Test whether the attenuation matches the formula in
        Theorem~\ref{thm:projection}(b).
\end{enumerate}

All four outcomes are informative: attenuation confirms the Projection Theorem;
no attenuation suggests either proxy noise, model misspecification, or that $\RD$
adds no explanatory power for this pair subset; partial attenuation provides an
estimate of $\operatorname{Var}(\widehat{R}) / (\operatorname{Var}(\widehat{R}) +
\operatorname{Var}(\zeta))$.

%=====================================================================
\section{Consequences for the Compression Account}
\label{sec:consequences}
%=====================================================================

\subsection{Family-level heterogeneity}

The Domain Divergence Theorem predicts that cross-family variation in partial
colexification should cluster in meaning pairs with high cross-community
variation in action-consequence structure---pairs from ecological, technological,
or ritual domains where communities differ most in what they do with the relevant
distinctions. This is testable: it predicts a correlation between the variance
of $\widehat{R}$ across communities and the variance of colexification pattern
across families for each meaning pair.

A concrete illustration is provided by the case of \textsc{spoon} and the
Chinese \textit{t\={a}ngch\'{\i}}.
Cross-linguistic databases treat these as translation equivalents mapping to the
same concept. Wierzbicka~\cite{wierzbicka2015} argues, however, that the two
words are embedded in different patterns of cultural reasoning and discourse,
owing to the different cultural functions of the distinct artefacts they
denote. In the reachability framework, this is not a methodological limitation
but a predicted phenomenon. The two words have similar $\RD$ in a
translation-equivalence domain $\theta_1$:

\[
\RD[\theta_1](\textsc{spoon},\; \textit{t\={a}ngch\'{i}}) \approx 0,
\]

because in that domain the relevant distinction between them does not control
different future trajectories. But in domains $\theta_2$ where eating practices,
culinary traditions, and social contexts of food use differ substantially between
communities, the reachability divergence increases:

\[
\RD[\theta_2](\textsc{spoon},\; \textit{t\={a}ngch\'{i}}) > 0.
\]

\begin{proposition}[Domain-Relative Lexical Divergence]
\label{prop:cultural}
Two words that are translation equivalents in a domain-general database may have
non-zero reachability divergence in culturally specific domains. The degree of
lexical differentiation between their language communities should be monotonically
related to $\RD[\theta_2]$ in those culturally specific domains: communities
with larger divergence in downstream object use should exhibit stronger lexical
differentiation than communities in which those downstream consequences remain
similar.
\end{proposition}

This converts Wierzbicka's cultural embedding observation from a methodological
limitation of cross-linguistic databases into a testable prediction: communities
exhibiting larger divergence in downstream utensil use should exhibit stronger
lexical differentiation, and English-derived similarity measures will be least
predictive precisely for meaning pairs with the highest cross-community variation
in $\RD$.

\subsection{The English-proxy distortion}

English-derived similarity and confusability measures will diverge most sharply
from the colexification patterns of other communities on meaning pairs where
English-speaking community future spaces differ from those other communities.
The reachability account predicts not merely that distortion exists but where
it is largest: in meaning pairs with the highest cross-community variation in
$\RD$.

\subsection{The interaction term}

Theorem~\ref{thm:projection} provides the structural interpretation of the
$\alpha_{SC}$ coefficient in (\ref{eq:standard}): it is the projection of
$\alpha_R \beta_{SC}$ from the omitted variable $\RD$ onto the $(S,C)$ plane.
The theorem converts the interaction from an empirical observation into a
derived quantity with a structural source and a testable consequence.

\subsection{Consequential ambiguity and reachability divergence}

A commentary on the Brochhagen et al.\ study by Beekhuizen~\cite{beekhuizen2026}
identifies a gap in the compression account that the reachability framework is
designed to fill. Beekhuizen notes that ambiguity is costly only when
disambiguation is consequential for the communicative goals at hand, and that
being maximally informative is not necessarily always one of those goals. He
treats this as a limitation of the distributional semantic measure used in the
paper, without being able to supply a principled criterion for when
disambiguation becomes necessary.

The reachability account supplies that criterion directly. Disambiguation becomes
necessary precisely when reachability divergence $\RD(X,Y)$ is sufficiently
large. When $\RD(X,Y) \approx 0$, the two meanings open the same admissible
futures in the relevant domain; merger is admissible and disambiguation provides
no navigational benefit. When $\RD(X,Y) \gg 0$, merger destroys navigational
structure and lexical distinction persists under repair pressure.

Beekhuizen's observation therefore converts from a criticism of the compression
account into a prediction of the reachability account: the cases where
distributional similarity fails to predict colexification patterns should be
exactly the cases where $\RD$ is high but $C$ is low. These are the meaning
pairs that are contextually non-confusable yet have divergent downstream
consequences---the (H, L, H) cell that the pilot table identifies as the primary
site of compression account misprediction.

%=====================================================================
\section{Language as Reachability Preservation}
\label{sec:language}
%=====================================================================

The five theorems together support a stronger thesis about the function of
language---though the thesis requires careful statement.

The compression account holds that language balances communicative efficiency
against ambiguity. This account is correct as far as it goes. Languages routinely
compress meaning, reuse forms, and exploit context to reduce lexical complexity,
and the empirical evidence for communicative efficiency as a driver of lexical
structure is strong. We do not dispute this.

What the reachability account adds is a constraint on compression: some
distinctions cannot remain compressed because doing so would foreclose futures
a community needs to keep open. This constraint sets the floor below which
compression cannot go without triggering repair. The compression account explains
why meanings merge; the reachability account explains why they sometimes refuse
to remain merged, and how they resist.

This motivates a revised thesis:

\medskip
\begin{quote}
\emph{Language is also a system for preserving the distinctions required to keep
a community's future navigable. Compression and differentiation determine which
meanings can be merged; reachability determines which merges can be sustained.}
\end{quote}
\medskip

The empirical signature of this constraint is exactly the phenomenon that
motivated this paper: distinctions that resist colexification despite high
semantic similarity and low contextual confusability. These are not anomalies
to be explained away by residual typological factors. They are the traces of
communities preserving futures they cannot afford to foreclose.

This reframes the relationship between language and knowledge. Lexical
distinctions are not merely representations of conceptual structure. They are
instruments of navigation. A language community that loses the distinction
between \textsc{hand} and \textsc{arm} in its medical vocabulary does not
merely become imprecise. It loses the representational capacity to route patients
to the correct clinical pathway. The distinction is preserved because its loss
is action-theoretically inadmissible---not because it is communicatively
ambiguous.

The same logic applies to legal distinctions (MURDER/MANSLAUGHTER,
THEFT/FRAUD), mathematical distinctions (CONTINUOUS/DIFFERENTIABLE,
CONVERGENT/CAUCHY), engineering distinctions (STRESS/STRAIN,
TORQUE/FORCE), and ecological distinctions (HABITAT/NICHE,
POPULATION/COMMUNITY). In each case the distinction is preserved not because
the terms are contextually confusable but because their collapse would remove
distinctions that control which futures are accessible.

%=====================================================================
\section{Scope, Limitations, and the Research Programme}
\label{sec:programme}
%=====================================================================

\subsection{What this paper establishes}

\begin{enumerate}
  \item \emph{Logical independence}: $\RD$ and $C$ are logically independent
        (Proposition~\ref{prop:independence}), demonstrated by colexification-relevant
        examples where both variables are in play.
  \item \emph{Sound principal results}: Three theorems and one proposition are
        derived from explicit definitions and assumptions rather than definitional
        stipulations. The Stability Theorem follows from graded repair dynamics
        (Assumptions~\ref{asm:repair} and~\ref{asm:convergence}). The Repair
        Theorem is an optimality result over admissibility-restoring repairs,
        with a scope remark bounding the uniqueness claim. The Domain Divergence
        Proposition is a direct consequence of domain-indexed future spaces. The
        Projection Theorem follows from omitted-variable algebra with explicitly
        stated confounding assumptions. The merge operator is governed by three
        named conditions whose joint implication ($\A_X \cap \A_Y \subseteq
        \A_{X \sqcup Y} \subseteq \A_X \cup \A_Y$) is stated and used in
        Lemma~\ref{lem:bounds}.
  \item \emph{Operationalizability}: $\RD$ can be estimated by at least three
        methods (Section~\ref{sec:operationalization}).
  \item \emph{Pilot estimation}: The pilot table (Section~\ref{sec:pilot}) shows
        that $\widehat{R}$ varies independently of $S$ and $C$ and predicts
        colexification patterns in the cases the compression account
        underdetermines.
  \item \emph{Falsifiable prediction}: The Projection Theorem generates a
        testable prediction about the interaction term in the Brochhagen et al.\
        model.
\end{enumerate}

\subsection{What this paper does not establish}

The framework is a theoretical research programme. The pilot table provides
directional evidence but not quantitative validation. The theorems are correct
given their assumptions; whether those assumptions hold in natural language
systems is an empirical question.

\subsection{Principal limitations}

\paragraph{Assumption sensitivity.}
The Stability Theorem depends on Assumptions~\ref{asm:repair}
and~\ref{asm:convergence}. Assumption~\ref{asm:convergence} idealises: real
languages may be slow to repair, may tolerate persistent inadmissibility in
peripheral domains, or may repair through routes that do not fully restore
navigational admissibility.

\paragraph{Domain boundary.}
The domain index $\theta$ is treated as given. In practice, domain boundaries
are fuzzy and historically variable. A fuller theory would need to explain how
$\theta$-partitions are maintained and how domain-boundary shifts affect
stability predictions.

\paragraph{Merge operator dependence.}
Lemma~\ref{lem:bounds} and the Stability Theorem hold for any admissible merge
operator. But the specific value of $\Lam(X,Y)$ depends on the choice of
$\sqcup$. Different choices yield the same qualitative predictions (stable vs.
unstable) but different quantitative values of $\widehat{R}$.

\paragraph{Pilot estimation scope.}
The $\widehat{R}$ estimates in Section~\ref{sec:pilot} are expert judgments, not
corpus-derived measurements. Full validation requires the methods in
Section~\ref{sec:operationalization}.

\subsection{Next steps}

\emph{Near term}: The empirical test of Theorem~\ref{thm:projection} against
Lexibank data with corpus-derived $\widehat{R}$ estimates for medical and legal
meaning pairs. This is the single most important validation step.

\emph{Medium term}: Diachronic testing of Theorem~\ref{thm:stability}(b) using
historical colexification data. The prediction that high-$\Lam$ colexifications
fragment faster is testable against diachronic corpora with domain annotations.

\emph{Longer term}: Cross-family testing of the Domain Divergence Theorem using
$\widehat{R}$ estimates from communities with different ecological, institutional,
and technological structures for the same meaning pairs.

%=====================================================================
\section{Conclusion}
\label{sec:conclusion}
%=====================================================================

We began with a question the compression account cannot answer: why do certain
distinctions survive?

The answer offered here is that they survive because losing them forecloses
futures a community needs to keep open. Compression determines whether meanings
can be merged. Reachability determines whether they can remain merged.

The formal development rests on three principal theorems, one structural
proposition, and one statistical theorem, each derived from explicit assumptions
rather than definitional stipulations. The Navigational Collapse Theorem
identifies a representational failure mode---indexed via an explicit measurable
performance function---that contextual confusability cannot detect. The Stability
Theorem, derived from graded repair dynamics rather than definitional circularity,
generates a diachronic prediction about the rate of colexification fragmentation.
The Repair Theorem, derived from a cost trade-off between representational
complexity and reachability loss, explains why partial colexification is the
preferred admissibility-restoring repair rather than full differentiation; the
scope of the uniqueness claim is precisely bounded to this class of repairs.
The Domain Divergence Proposition---a direct structural consequence of domain
indexing rather than a non-trivial derivation---predicts the systematic emergence
of technical vocabularies and explains cross-family heterogeneity in partial
colexification. The Projection Theorem generates a directly falsifiable prediction
about the interaction term in the best existing statistical model of
colexification.

Before these theorems, we established that reachability divergence and contextual
confusability are logically independent quantities, demonstrated by
colexification-relevant meaning pairs where the critical dissociation---low $C$,
high $\RD$---occurs among pairs that are semantically similar and cross-linguistically
attested as colexification candidates. We then presented a pilot estimation showing
that $\widehat{R}$ varies independently of $S$ and $C$ and predicts colexification
patterns in the cases the compression account underdetermines.

The compression account explains a great deal about lexical structure. It explains
why meanings merge, which merges are stable in the absence of domain-specific
pressure, and why partial colexification emerges as a compromise between
efficiency and disambiguation. The reachability account does not replace it. It
adds a constraint the compression account cannot express: the floor below which
compression cannot go without triggering repair. At that floor, lexical
distinctions survive not because they are communicatively necessary, but because
their loss is action-theoretically inadmissible.

Together these results suggest that the study of lexical organisation is also,
at its deepest level, the study of how communities preserve the futures they need
to remain navigable.

Bloomfield famously characterised the lexicon as a collection of basic
irregularities~\cite{bloomfield1933}. The reachability account proposes the
opposite interpretation. Lexical distinctions appear irregular only when viewed
through compression alone. Once reachability divergence is taken into account,
many apparently arbitrary distinctions emerge as systematic responses to
differences in admissible futures: the distinctions that survive are exactly
those whose loss would be action-theoretically inadmissible. The irregularity
is in the compression account's model of language, not in language itself.

%=====================================================================
\begin{thebibliography}{99}

\bibitem{beekhuizen2026}
Beekhuizen, B.\ (2026).
When languages favour complex words.
\textit{Nature Human Behaviour}.
\url{https://doi.org/10.1038/s41562-026-02493-6}

\bibitem{bloomfield1933}
Bloomfield, L.\ (1933).
\textit{Language}.
Henry Holt.

\bibitem{brochhagen2022}
Brochhagen, T.\ \& Boleda, G.\ (2022).
When do languages use the same word for different meanings?
The Goldilocks principle in colexification.
\textit{Cognition}, 226, 105179.

\bibitem{brochhagen2026}
Brochhagen, T., Liao, X., Wright, J.\ D., \& Saldana, C.\ (2026).
The interaction of meaning similarity and confusability explains regularity
in form--meaning mappings at and below the word level.
\textit{Nature Human Behaviour}.
\url{https://doi.org/10.1038/s41562-026-02488-3}

\bibitem{francois2008}
Fran\c{c}ois, A.\ (2008).
Semantic maps and the typology of colexification:
Intertwining polysemous networks across languages.
In M.\ Vanhove (Ed.),
\textit{From Polysemy to Semantic Change}, 163--215.
John Benjamins.

\bibitem{gardenfors2000}
G\"{a}rdenfors, P.\ (2000).
\textit{Conceptual Spaces: The Geometry of Thought}.
MIT Press.

\bibitem{karjus2021}
Karjus, A., Blythe, R.\ A., Kirby, S., Wang, T., \& Smith, K.\ (2021).
Conceptual similarity and communicative need shape colexification:
an experimental study.
\textit{Cognitive Science}, 45, e13035.

\bibitem{kemp2012}
Kemp, C.\ \& Regier, T.\ (2012).
Kinship categories across languages reflect general communicative principles.
\textit{Science}, 336, 1049--1054.

\bibitem{list2022}
List, J.-M.\ et al.\ (2022).
Lexibank, a public repository of standardized wordlists with computed
phonological and lexical features.
\textit{Scientific Data}, 9, 316.

\bibitem{piantadosi2012}
Piantadosi, S.\ T., Tily, H., \& Gibson, E.\ (2012).
The communicative function of ambiguity in language.
\textit{Cognition}, 122, 280--291.

\bibitem{regier2016}
Regier, T., Carstensen, A., \& Kemp, C.\ (2016).
Languages support efficient communication about the environment:
words for snow revisited.
\textit{PLOS ONE}, 11(4), e0151138.

\bibitem{xu2020}
Xu, Y., Duong, K., Malt, B.\ C., Jiang, S., \& Srinivasan, M.\ (2020).
Conceptual relations predict colexification across languages.
\textit{Cognition}, 201, 104280.

\bibitem{wierzbicka2015}
Wierzbicka, A.\ (2015).
Language and cultural scripts.
\textit{Language Sciences}, 47, 66--83.

\bibitem{zaslavsky2018}
Zaslavsky, N., Kemp, C., Regier, T., \& Tishby, N.\ (2018).
Efficient compression in color naming and its evolution.
\textit{Proceedings of the National Academy of Sciences}, 115, 7937--7942.

\end{thebibliography}

\end{document}
